% Generated draft: requires mathematical and source review.
\documentclass[11pt]{article}
\usepackage[margin=25mm]{geometry}
\usepackage{fontspec}
\setmainfont{Latin Modern Roman}
\usepackage{amsmath,amssymb,mathtools,unicode-math}
\setmathfont{Latin Modern Math}
\usepackage{enumitem,needspace,xurl,hyperref}
\setlength{\parindent}{0pt}
\setlength{\parskip}{5pt}
\setlength{\emergencystretch}{3em}
\newcommand{\N}{\mathbb{N}}
\newcommand{\Z}{\mathbb{Z}}
\newcommand{\Q}{\mathbb{Q}}
\newcommand{\R}{\mathbb{R}}
\newcommand{\C}{\mathbb{C}}
\newcommand{\F}{\mathbb{F}}
\newcommand{\A}{\mathbb{A}}
\newcommand{\PP}{\mathbb{P}}
\newcommand{\E}{\mathbb{E}}
\newcommand{\eps}{\varepsilon}
\newcommand{\dd}{\,\mathrm{d}}
\begin{document}
\section*{UnsolvedMath ID: 800912}
\textbf{Source problem code:} catalogue entry
\subsection{Definitions and mathematical statement}
Let \(X\) be a compact K\"ahler manifold of complex dimension \(n\). The supplied question concerns a positive current \(T\) that is harmonic in the specific sense
\[
 dd^cT=0.
\]
Here one may fix the convention \(dd^c=\frac{i}{\pi}\partial\bar\partial\). The assumption is \(dd^c\)-closedness; the source does not assume that \(T\) is \(d\)-closed.

For a current of pure bidimension \((p,p)\), its Lelong number at \(z\), in local holomorphic coordinates \(w\) centered at \(z\), has the usual normalized mass-density expression
\[
 \nu(T,z)=\lim_{r\downarrow0}\frac{1}{2^p r^{2p}}
 \int_{\{\lVert w\rVert<r\}}T\wedge\bigl(dd^c\lVert w\rVert^2\bigr)^p.
\]
The source states that Lelong numbers are well-defined for positive currents satisfying \(dd^cT=0\). Its notation \(v(T,z)\) is written here as \(\nu(T,z)\). The source does not specify the bidimension of \(T\), its permitted range, or whether currents with mixed degrees are intended; the displayed definition describes the usual pure-degree setting, and the intended degree scope requires confirmation.

\textbf{Question 44 (attributed to Dinh).} For a compact K\"ahler manifold \(X\) and a positive harmonic current \(T\) as above, are all the upper Lelong level sets
\[
 E_c(T)=\{z\in X:\nu(T,z)\geq c\}
\]
complex analytic subsets of \(X\)? The source leaves the range of \(c\) unstated. For real thresholds, the substantive question is whether this holds for every \(c>0\); nonnegativity of Lelong numbers gives \(E_c(T)=X\) for \(c\leq0\). No codimension or purity requirement for these analytic subsets is stated.
\subsection{Short English statement}
The question attributed to Dinh asks whether every positive \(dd^c\)-closed current on a compact K\"ahler manifold has complex analytic upper Lelong level sets. The source does not specify the current's degree or the threshold range. The registry labels the problem open; this is supplied status information, not an independently verified assessment.
\subsection{Sources}
\begin{itemize}
\item Local UnsolvedMath registry, permanent ID \(800912\), exact source problem code \texttt{catalogue entry}; supplied title: \emph{Open Problems in Pluripotential Theory --- Question 44 (Dinh)}. The registry lists Dinew, Guedj, and Zeriahi as proposers, with year \(2015\), and status \texttt{open}.
\item Dinew, Guedj, and Zeriahi, \emph{Open Problems in Pluripotential Theory} (2015), Question 44; attribution/context: Dinh. Supplied reference: \url{https://arxiv.org/abs/1511.00705}.
\item The supplied provenance records access on \texttt{2026-10-16} and extraction by \texttt{arxiv\_latex}. These are supplied provenance claims, not access or source verification performed for this draft. No separate dated literature assessment is supplied.
\end{itemize}
\end{document}
